% 本文件由 main.tex 输入，作为“文献综述与理论假说”章节，不单独编译。

\section{文献综述与理论假说}

\subsection{文献综述}

\subsubsection{基金促进创新的效应与机制的已有研究}

围绕政府产业引导基金能否促进创新，现有研究大体形成了“促进论”与“约束论”两类观点。较多研究认为，引导基金本质上是财政资金“基金化”和产业政策“金融化”的重要载体，其通过股权投资而非传统补贴方式介入企业创新活动，有助于克服科技创新中的融资约束、信息不对称和风险外部性，从而提升企业创新投入与创新产出。吴超鹏和严泽浩（2023）发现，政府基金引导能够显著促进企业在关键核心技术领域的创新，其重要原因在于政府资本提高了对创新失败的容忍度，并缓解了企业外部融资约束。刘玉斌等（2023）基于战略性新兴产业创业投资引导基金的研究也表明，政府引导基金对企业创新具有显著促进作用，其机制主要体现为融资约束缓解、认证效应与信号传递。蔡庆丰等（2024）则进一步从区域层面指出，产业引导基金不仅能够直接促进域内企业创新投入，而且能够通过资本集聚形成融资效应、信息效应和竞争效应，放大财政资金“拨改投”的微观创新激励。

不过，另一部分研究对引导基金的创新绩效持更为审慎的判断。这类研究强调，政府资本虽然具有政策目标，但也可能因为多重委托代理、风险规避、考核约束和寻租行为而偏离“投早、投小、投科技”的初衷。徐明（2022）发现，政府风险投资不仅未必促进企业创新，反而可能因目标冲突、治理失灵和政治关联而抑制企业研发投入和创新产出。进一步地，部分文献指出，引导基金在实际运作中可能偏好相对成熟、风险较低的项目，由此形成对民间风险资本的挤出或同质化竞争，而非真正填补创新融资缺口。由此可见，现有研究并不否认引导基金具有促进创新的潜在能力，但其创新绩效能否实现，很大程度上取决于基金是否能够保持耐心资本属性、有效撬动社会资本，并真正流向高风险、长周期的创新活动。

总体来看，关于“基金是否促进创新”的争论，已经从简单的政策有效性判断，转向对其作用条件和内在机制的识别。这为本文进一步讨论债务约束条件下政府产业引导基金何以促进创新、又何以可能偏离创新目标，提供了直接的文献基础。

\subsubsection{地方财政债务压力对创新的已有研究}

与引导基金研究相对应，另一条重要文献线索聚焦于地方财政压力和政府债务扩张对创新活动的影响。既有研究普遍认为，财政压力与债务约束并非单纯的宏观背景变量，而是会通过改变政府行为、金融资源配置和企业融资环境，对创新形成系统性影响。李森和王聪（2024）发现，财政压力会显著抑制制造业企业创新，并通过加剧政治晋升激励、减少创新补贴和提升企业税负等渠道发挥作用。类似地，部分研究指出，在财政收支矛盾加剧时，地方政府更可能强化税收征管、压缩长期性科技支持支出，从而弱化企业开展高风险研发活动的能力和意愿。

从债务扩张的角度看，现有研究更强调其金融挤出效应与资源错配后果。余海跃和康书隆（2020）发现，地方政府债务扩张会通过抬升企业融资成本挤出企业投资，且对高效率民营企业的影响更强。徐彦坤（2020）的研究进一步表明，地方政府债务上升会加剧企业融资约束、抬高债务资本成本并压缩投资机会。刘贯春等（2022）则从投融资期限结构角度指出，地方政府债务治理有助于缓解政企融资竞争并改善企业“短贷长投”问题，其反向含义是，债务扩张本身会恶化长期资金供给。除融资挤出外，地方政府债务还可能改变企业的资产配置方向。孔丹凤和谢国梁（2021）发现，地方政府债务扩张会加剧企业金融化倾向，推动资金从实体创新活动转向短期金融收益。

更进一步，财政压力和债务压力还会改变地方政府干预市场的方式。张跃文等（2025）指出，在财政压力上升背景下，地方政府可能通过“合作型”干预影响企业投资决策，以增加短期税源和财政收入。这说明，财政压力不仅影响政府“有没有钱支持创新”，还会影响政府“如何使用政策工具”。如果把这一逻辑延伸到政府产业引导基金，就意味着地方债务压力可能促使地方政府更加重视资金安全、回收速度和短期绩效，进而削弱基金原本应有的耐心资本属性、让利属性和创新导向属性。

因此，现有关于财政债务压力的研究已经揭示出一个重要事实：地方政府面临的债务约束会同时影响企业创新环境和政府自身的行为逻辑。只是现有文献更多考察财政压力或地方债务对企业创新的直接影响，对于其如何通过扭曲政府产业引导基金这一政策工具而间接影响创新，仍缺乏专门讨论。

\subsubsection{现有文献的不足与本文边际贡献}

综合来看，现有研究虽然分别从政府产业引导基金和地方财政债务压力两个维度积累了较为丰富的成果，但仍存在三方面不足。第一，两条文献线索大多彼此分离。已有研究要么关注引导基金能否促进企业创新，要么考察财政压力和地方债务如何抑制企业创新，较少将“地方债务压力—引导基金运作—企业创新绩效”纳入统一分析框架，从而难以解释为何同样是政府产业引导基金，在不同财政约束情境下可能表现出显著差异。第二，既有研究对于引导基金失灵的讨论，虽然涉及风险规避、代理问题和资本挤出，但较少进一步从耐心资本属性、社会资本撬动效率和融资约束缓解效应三条机制链进行系统展开，因而对债务压力如何改变基金行为逻辑的解释仍不够细化。第三，现有研究更多停留在“基金是否有效”的平均效应判断上，对政策工具的作用边界讨论不足，尤其缺少对地方债务压力这一现实制度约束的关注。

基于上述不足，本文尝试将地方财政债务压力、政府产业引导基金与创新绩效纳入统一分析框架，重点考察债务压力如何削弱政府产业引导基金促进创新的作用，并从耐心资本属性扭曲、社会资本撬动效率下降和融资约束缓解效应弱化三个方面展开理论分析。相较于既有文献，本文的边际贡献主要体现在：其一，不再停留于讨论“引导基金是否促进创新”的一般性命题，而是进一步回答在债务约束背景下这一促进作用为何会减弱；其二，将宏观债务约束与中观政策工具运作结合起来，拓展了地方债务压力影响创新的传导链条研究；其三，从政策工具作用边界的角度，为理解政府产业引导基金何以“有效”、又何以“失效”提供新的分析视角。基于此，下文将进一步提出本文的理论假说。

\subsection{理论分析和理论假说}

企业的创新活动因其长期性和不确定性，受到严重的融资约束，主要依靠企业内部资金( Hall，2002)。但企业利用内部资金进行创新活动存在两大难题：一是内部财务波动会直接影响创新投入，易受到外部冲击，可能导致资金链断裂，创新活动被迫中断；二是创新活动具有较高的调整成本，突然中断和再重新启动，会使企业承担巨大损失（Hall，2002，2005）。为解决这一问题，受融资约束的企业会根据调整成本的高低进行反向调整，优先削减调整成本低的投资，将有限资金保留用于调整成本高的创新活动，最后保持各种投资的净收益在边际上相等，从而保证创新投资的连续性与稳定性( Fazzari et al.1993)。

然而，初创期和早中期的科技型企业普遍存在企业创新资金不足、外部融资成本过高、融资渠道不通畅的问题。从市场失灵理论来看，战略性新兴产业企业的创新活动面临三重典型的市场失灵：一是信息不对称，企业技术前景不确定性高、投资周期长的特点，外部投资者需付出较高成本以准确评估其偿债能力与成长潜力；二是正外部性，企业创新成果存在知识溢出效应，社会收益高于私人收益，市场化资本在收益风险权衡下天然倾向于投资不足；三是融资缺口，企业初创期和早中期缺乏稳定的现金流与可抵押资产。在这一背景下，企业内部可以用来平滑创新波动的资本十分有限，单靠市场机制作用与企业自身管理难以有效缓解其面临的较强融资约束。

为应对这一问题，政府产业引导基金作为一项纠正市场失灵的政策性金融安排，为战略性新兴产业企业突破融资约束提供了可行路径。政府产业引导基金参股基金根据企业所处生命周期的不同阶段，实施阶段性投资与联合投资，匹配企业差异化的资金需求，直接为企业创新活动提供资本支持。同时，政府产业引导基金的发展有利于促进地区金融资源的均衡分布，资本资源聚集的网络效应有利于企业信息在资本之间的共享，降低社会资本寻求投资对象过程中的信息成本（李健旋和赵林度，2018）。通过上述机制，政府产业引导基金不仅直接补充了企业创新资金，更改善了企业融资条件，为社会资本进入创新领域提供保障，同时有助于企业保持创新投资的稳定性与连续性。

由此提出假说H1：政府产业引导基金能够促进城市创新。

然而，引导基金的杠杆运作模式在放大财政资金效果的同时，近年来也面临地方债务压力日益严峻的现实问题。截至2024年末，地方政府法定债务余额已达47.54万亿元，隐性债务约14.3万亿元，地方整体偿债压力持续攀升（国务院，2025）。Wang等（2025）在《Sustainability》上的研究发现，地方政府债务扩张显著降低企业创新质量，且财政压力较低的地区这一负面影响更小。Zhang和Jin（2022）发现地方债务与创新的倒U型关系，拐点为债务与GDP之比0.689。这些发现告诉我们，债务压力可能改变了引导基金支持创新产出的制度环境和运作条件。债务压力可能会通过恶化融资约束，侵蚀耐心资本属性，减少社会资本撬动作用等途径，影响产业引导基金支持创新的效果和能力。

由此提出假说H2:地方债务压力会削弱政府产业引导基金的创新促债务压力调节效应。

创新活动的长周期性、高不确定性和正外部性这三个特征共同决定了创新活动对耐心资本的内生需求：一种能够长期持有、耐受不确定性、不追求短期回报退出的资本形态。理解耐心资本作为引导基金的本质，需要区分两个层次。第一个层次是功能层次：引导基金填补市场耐心资本的供给缺口，在市场资本“不愿投、不敢投”的领域发挥补位作用。第二个层次是信号层次：引导基金的存在本身就是对创新长期价值的制度性背书，其“耐心”属性向市场传递了一个关键信号——某些创新项目的价值需要以十年为单位衡量，而非以季度或年度衡量。这个信号具有公共品性质，私人资本不会自发提供。

引导基金对地方创新产出的支持效果，从根本上取决于其耐心资本属性能否完整保持, 将资金投向国务院要求的“投早投小投科技”的方向。一旦这一属性被侵蚀，引导基金就不再能够履行其制度功能，其支持创新产出的效果必然下降。

首先，偿债支出具有预算优先权，在财政资源有限时优先于引导基金出资。这可能会导致新设基金被推迟或取消，已有基金的后续出资被延缓，并且基金存续期被迫缩短；其次，基金倾向于追求保值增值导致偏向成熟项目，偏离“投早投小”的目标。引导基金的政策初衷是支持处于早期阶段、创新含量高但风险也高的科技企业。在债务压力下，地方政府对财政资金保值增值的要求更为严格，国有资产考核体系使得管理人“不敢投”初创期企业，宁愿投向确定性高但支持必要性不高的成熟期项目，甚至购买理财产品。

考核体系的变形使得地方政府倾向于缩短考核周期，要求基金尽快呈现正向回报。田轩指出，\erjfootnote{北京大学博雅特聘教授田轩在2026博鳌亚洲论坛年会上的相关发言。}国资考核机制与耐心资本的市场化表现存在矛盾——创投容错机制不健全，收益负面波动时缺乏包容甚至问责处罚。当考核周期被压缩，基金管理人不再有“耐心等待”创新产出的空间，投资决策从“选择最有长期价值的创新项目”转向“选择能在考核期内呈现正回报的安全项目”。审计发现，山西8.12亿元基金违规投向房地产开发等项目，\erjfootnote{山西省审计厅《关于2023年度省级预算执行和其他财政收支的审计工作报告》。}云南某公司借用政府投资基金1.54亿元用于偿还债务。\erjfootnote{云南省人大常委会2023年度审计查出问题整改情况报告（2024年11月发布）。}当债务压力较低时，地方政府能够容忍引导基金的合理亏损，投资决策向“投早投小”倾斜；当债务压力较高时，投资决策向“保本保息”倾斜，创新最需要支持的早期科技企业反而得不到资金。Zhao和Zhang（2025）的研究也证实了这一点——引导基金在高新技术产业和创投成熟的地区对创新边界的拓展效应更强，而这些地区恰恰是债务压力相对较低的地区。最后，值得考虑的是，在债务压力下，行政干预加剧会造成投资决策变形，进而侵蚀耐心资本的属性。面临债务压力的地方政府，更可能将引导基金视为缓解财政紧张或实现短期政绩的工具，加强行政干预。审计显示，宁夏4400万元基金因要求的投资收益率过高导致3520万元闲置，\erjfootnote{宁夏2023年度重大政策落实及重大项目审计情况（2024年8月发布）。}辽宁1只基金到位5.02亿元仅投资4830万元（占比9.66\%）。\erjfootnote{辽宁省审计厅相关审计报告（2024年发布）。}Wang等（2025）发现政府干预较少、财政压力较低的地区，债务对创新质量的负面影响更小，耐心资本属性发挥的作用会更强。

由此提出假说H3：债务压力通过削弱引导基金的耐心资本属性，进一步抑制其创新促进作用。

政府产业引导基金能以少量财政资金撬动数倍社会资本，形成规模化创新资本供给（Xiao et al., 2024；蔡庆丰等，2024）。然而，地方政府债务压力从多个维度削弱了这一撬动效率，抑制了企业创新。

首先，债务压力会挤占财政空间，导致引导基金出资额减少（崔竹，2025）。当地方政府债务压力较低时，财政资源相对充裕，政府能够按计划足额实缴承诺出资，引导基金得以顺利运作。当债务压力较高时，偿债支出和“三保”等刚性支出挤占财政空间，政府可用于引导基金的财政资金被严重压缩，社会资本的募集减少；其次是债务压力会诱致“名股实债”等异化行为，即引导基金从风险投资工具异化为债权工具。为弥补社会资本信心不足，债务压力大的地区更倾向于通过承诺回购本金、设定最低投资收益或承担投资损失等方式吸引社会资本（崔军等，2025）。社会资本的“撬动”功能被“替代”功能所取代。罗志恒（2021）测算，引导基金通过杠杆承诺、劣后级兜底等机制使地方政府隐性债务增加约1.96至4.54万亿元，形成“债务越高—越需要保底募资—隐性债务越膨胀—创新功能越弱”的恶性循环；此外，社会资本对高债务地区信心下降，募资难度加大，直接降低了引导基金对社会资本的撬动效率（张圆圆、夏球，2024）。当债务压力较低时，地方财政实力强、信用好，社会资本参与意愿高；当债务压力较高时，社会资本对地方财政信心下降，社会资本在高债务地区的投资意愿显著下降，企业创新获得的总资本量大幅减少，撬动效率被严重削弱。撬动效率的降低直接限制了可用于支持企业创新的资金规模，企业创新的融资约束难以得到有效缓解（Hao et al., 2025）。

由此提出假说H4:债务压力通过降低社会资本撬动效率，进一步抑制基金创新促进作用。

初创期和早中期的科技型企业面临三重典型的市场失灵：信息不对称、正外部性、缺乏用于融资的可抵押资产，三重作用叠加使得其难以仅靠市场机制作用与企业自身管理缓解融资约束。相较于依赖稳定现金流和抵押资产的债权融资，股权融资更能分担科技创新活动中的高风险与长周期特征，也更契合科技型企业持续研发投入的资金需求。政府产业引导基金根据企业所处生命周期的不同阶段，实施阶段性投资与联合投资，匹配企业差异化的资金需求，直接为企业创新活动提供长期资本支持。同时，政府产业引导基金的发展有利于促进地区金融资源的均衡分布，资本资源聚集的网络效应有利于企业信息在资本之间的共享，降低社会资本寻求投资对象过程中的信息成本（李健旋和赵林度，2018）。因此，政府产业引导基金不仅能够通过股权投资直接缓解科技型企业的资金缺口，还能够通过政府背书、风险分担与联合投资向社会资本释放质量认证信号，进而提升企业获得外部股权融资的可能性。

然而地方债务压力的攀升从源头上削弱产业引导基金缓解企业融资约束的能力。地方政府债务过高源于政府投资性项目偏离实际公共需求，从而造成政府支出低效率配置，不利于当地经济增长（陈思霞和陈志勇，2015)。在此背景下，债务率较高地区的地方财政逐渐进入“紧平衡”状态。在经济下行压力下，地方政府面临融资渠道受限与债务风险上升等问题（崔竹，2025），财政空间进一步收窄，从而压缩了政府对产业引导基金持续出资的能力。我国各级政府投资基金的累计认缴率约为64.9\%，大量已认缴资金无法实际到位，直接影响引导基金的正常运作。更重要的是，地方债务压力会削弱政府产业引导基金作为股权融资工具的风险承担能力与资本撬动能力。一方面，财政紧平衡会使地方政府更倾向于将有限资金投向短期稳增长和债务化解任务，降低对初创期、早中期科技型企业这类高风险、长周期项目的持续支持力度；另一方面，财政出资不到位和政策连续性下降会削弱政府资金的信用背书功能，降低社会资本对引导基金的跟投意愿，使基金原本通过联合投资放大资本供给、缓解信息不对称的机制难以充分发挥。除了直接影响资金来源外，地方债务压力还会通过压缩整体融资环境进一步削弱政府投资基金的作用。研究指出，地方政府债务扩张会通过直接占压银行信贷资源，对企业融资形成“挤出效应”，这种挤出效应在“中小企业和非国有企业”中表现得比大型企业和国有企业更为显著（华夏等，2020）。我国金融市场银行信贷的配置效率较低下，在缺乏融资的抵押资产初创企业中，银行更加青睐有政府信用作为隐性担保的国有企业和融资平台（周天宇，2025）。当引导基金自身出资能力受限、社会资本撬动效率下降，而外部融资环境又因债务扩张而趋于收紧时，初创期和早中期科技型企业面临更强的股权融资约束，其融资约束被进一步恶化，从而严重制约研发投入与持续创新能力的提升。

由此提出假说H5：债务压力通过弱化基金缓解融资约束的能力，进一步抑制城市创新。